\chapter{Tecnologías y Herramientas Utilizadas}
\label{cap:tecnologias_utilizadas}

El desarrollo de la plataforma \textit{SkiClubManager} se ha basado en una arquitectura desacoplada compuesta por un servicio de \textit{backend} centralizado y una aplicación móvil nativa como \textit{frontend}. La selección de la pila tecnológica (\textit{stack}) responde a criterios estrictos de rendimiento, mantenibilidad y la necesidad de ofrecer una experiencia de usuario fluida en entornos de alta montaña. A continuación, se justifican técnicamente las herramientas y entornos seleccionados.

\subsubsection{Infraestructura Backend y API REST}
Para dar soporte a la aplicación móvil y centralizar la lógica de negocio, se ha diseñado un servicio backend basado en una arquitectura cliente-servidor mediante el desarrollo de una API REST.

\subsubsection{Node.js}
Node.js es un entorno de ejecución para JavaScript construido con el motor V8 de Google \cite{NodeJS}. Su elección como núcleo del \textit{backend} se justifica por las siguientes características:
\begin{itemize}
    \item \textbf{Modelo de E/S sin bloqueo y orientado a eventos:} Utiliza un único hilo de ejecución asistido por un bucle de eventos (\textit{Event Loop}). Esto permite gestionar múltiples conexiones concurrentes de forma asíncrona sin bloquear el servidor, una cualidad crítica cuando entrenadores y familias interactúan con la aplicación simultáneamente.
    \item \textbf{Unificación del lenguaje:} El uso de JavaScript tanto en el servidor como en el cliente (\textit{frontend}) optimiza los tiempos de desarrollo, reduce la curva de aprendizaje y simplifica el mantenimiento del código fuente global.
\end{itemize}

\subsubsection{Express}
Express es un \textit{framework} minimalista y flexible creado para Node.js \cite{ExpressJS}, seleccionado para estructurar la API REST debido a su solvencia en el manejo de peticiones HTTP:
\begin{itemize}
    \item \textbf{Arquitectura basada en Middleware:} Organiza el flujo de las peticiones mediante funciones interconectadas. Esto permite auditar, transformar o validar los datos de entrada antes de que alcancen la lógica de negocio final.
    \item \textbf{Enrutamiento modular:} Permite desglosar las rutas de la API en controladores independientes, garantizando un código limpio, escalable y alineado con los principios de diseño de software modular.
\end{itemize}

\subsubsection{Motor de Persistencia: MySQL}
Como sistema de gestión de bases de datos se ha seleccionado MySQL \cite{MySQL}, un motor relacional de código abierto que garantiza el almacenamiento seguro y consistente de la información del club. Su elección se fundamenta en su capacidad para gestionar restricciones de integridad referencial mediante claves foráneas, asegurando que las asistencias, transportes y expedientes queden vinculados de forma unívoca y protegida a las fichas de los deportistas menores de edad.

\subsubsection{Desarrollo Frontend Móvil}
La capa de interacción con el usuario final requiere un entorno capaz de ofrecer portabilidad sin penalizar el rendimiento en dispositivos móviles.

\subsubsection{React Native}
React Native es un \textit{framework} de código abierto desarrollado por Meta \cite{ReactNative} que permite compilar aplicaciones móviles nativas para iOS y Android a partir de una única base de código en JavaScript y React. Los motivos técnicos de su elección incluyen:
\begin{itemize}
    \item \textbf{Renderizado Nativo:} A diferencia de las soluciones híbridas basadas en vistas web, React Native invoca directamente los componentes nativos del sistema operativo, garantizando una interfaz fluida y una respuesta inmediata a las interacciones del usuario.
    \item \textbf{Eficiencia en Entornos de Nieve:} La ligereza de la aplicación y la velocidad de respuesta de sus componentes son críticas para el cuerpo técnico en las pistas de esquí, donde las bajas temperaturas exteriores exigen flujos de trabajo rápidos y pantallas de carga mínimas.
\end{itemize}

\subsubsection{Gestión de Estado: Context API}
Para la propagación del estado global de la aplicación se ha utilizado la herramienta nativa \textit{Context API} de React. Este mecanismo permite inyectar los datos de la sesión activa, el rol del usuario autenticado (Entrenador o Padre) y los filtros temporales a lo largo de todo el árbol de componentes sin necesidad de recurrir a librerías externas complejas, facilitando el renderizado condicional inmediato de la interfaz.

\subsubsection{Autenticación: JSON Web Tokens (JWT)}
La seguridad y el aislamiento de los datos sensibles se gestionan mediante el estándar abierto RFC 7519, conocido como JSON Web Tokens \cite{RFC7519_JWT}. Tras un inicio de sesión exitoso, el servidor genera un token firmado criptográficamente que la aplicación móvil almacena de forma segura. Este token se adjunta en las cabeceras de cada petición HTTP, garantizando la autenticidad del emisor y permitiendo la validación de roles en cada \textit{endpoint} de la API REST.

\subsubsection{Visualización Analítica: React Native Chart Kit}
Para la representación gráfica del rendimiento de los corredores y los balances de asistencia en los perfiles familiares, se ha implementado la librería complementaria \textit{react-native-chart-kit} \cite{ReactNativeChartKit}. Esta herramienta se acopla directamente sobre el flujo de datos del estado global, transformando registros relacionales planos de asistencia en métricas visuales comprensibles y optimizadas para su correcta renderización en dispositivos móviles.

\subsubsection{Diseño de Interfaces: Figma}
Como entorno de prototipado industrial se ha empleado Figma. Esta herramienta basada en la nube ha permitido modelar la arquitectura de la información y evolucionar los esquemas desde \textit{wireframes} de baja fidelidad hasta los \textit{mockups} de alta fidelidad que sirven como plano técnico exacto para la maquetación visual de la aplicación móvil.